\part*{Anhang}
\addcontentsline{toc}{chapter}{Anhang}

\chapter{Allgemeine Daten}
\label{ch:a_grund}

\begin{table}[h!]
\centering
\caption{Zusammensetzung der trockenen Luft, nach \cite{hausen_tieftemperaturtechnik_1985, windmeier_cryogenic_2013}}
\label{tab:a_Luftzusammensetzung}
\begin{tabular}{l r}
\toprule
Komponente & Anteil \\
\midrule
Stickstoff (N$_2$)     & 78{,}084\,Vol.\% \\
Sauerstoff (O$_2$)     & 20{,}942\,Vol.\% \\
Argon (Ar)             & 0{,}934\,Vol.\% \\
Kohlendioxid (CO$_2$)  & 0{,}035\,Vol.\% \\
\midrule
Neon (Ne)              & 18{,}2\,ppm \\
Helium (He)            & 5{,}2\,ppm \\
Krypton (Kr)           & 1{,}1\,ppm \\
Xenon (Xe)             & 0{,}09\,ppm \\
Wasserstoff (H$_2$)    & 0{,}05\,ppm \\
Ozon (O$_3$)           & 0{,}003–0{,}03\,ppm \\
\midrule
\multicolumn{2}{l}{Wasserdampf (H$_2$O): variabel, typ.\ 25-75\,Vol.\% bei 15\,°C} \\
\bottomrule
\end{tabular}
\end{table}

\chapter{Modellierung}
\label{ch:a_sim}

\begin{landscape}
\centering
\begin{longtable}{lrrrrrrrrrrrr}
\caption{Thermodynamische und exergetische Kenngrößen der simulierten Prozessströme des Doppelkolonnenmodells} \\
\hline
Stream & $\dot m$ & $\dot n$ & $T$ & $p$ & $h$ & $s$ & $l_{frac}$ & $v_{frac}$ & $e_{PH}$ & $e_{CH}$ & $e_T$ & $e_M$ \\
 & (kg / s) & (mol / s) & (K) & (Pa) & (J / kg) & (J / kgK) &  &  & (J / kg) & (J / kg) &  &  \\
\hline
S1 & 30,725504099999998 & 1064,8131900000001 & 288,14999999999998 & 101325 & -100458,29300000001 & 127,232974 & 0 & 1 & 0 & 3840,6965306714042 & 0 & 0 \\
S2 & 30,725504099999998 & 1064,8131900000001 & 394,875269 & 262471,45500000002 & 7580,6939300000004 & 171,97623899999999 & 0 & 1 & 95146,214704415557 & 3840,6965306714042 & 16361,47017682291 & 78784,744495046383 \\
S3 & 30,725504099999998 & 1064,8131900000001 & 308,14999999999998 & 252471,45499999999 & -80634,300900000002 & -69,049189999999996 & 0 & 1 & 76382,697005124166 & 3840,6965306714042 & 802,13803554812966 & 75580,559148580403 \\
S4 & 30,725504099999998 & 1064,8131900000001 & 422,07465100000002 & 654000 & 34889,142099999997 & -24,2954632 & 0 & 1 & 179010,35348676349 & 3840,6965306714042 & 25139,359031704222 & 153870,994422513 \\
S5 & 30,725504099999998 & 1064,8131900000001 & 308,14999999999998 & 644000 & -85507,959400000007 & -354,02290499999998 & 0,0024337600800000002 & 0,99756624000000005 & 153624,21474477943 & 3840,6965306714042 & 1017,1238752759797 & 152607,09066771666 \\
S6 & 30,6800067 & 1062,2877100000001 & 308,02945 & 634000 & -62015,964099999997 & -336,44734499999998 & 0 & 1 & 152470,02635106986 & 3993,9181304024942 & 1005,6793957610186 & 151464,34697486553 \\
S7 & 0,0454973797 & 2,5254784099999998 & 308,02945 & 634000 & -15926740,5 & -9176,8422100000007 & 1 & 0 & 3588,5836739736465 & 49956,505736966654 & 2965,9820167621647 & 622,60166160733866 \\
S8 & 30,513200099999999 & 1053,6369699999998 & 308,00322499999999 & 624000 & 8290,0620400000007 & -331,74889899999999 & 0 & 1 & 150768,73533169666 & 5249,6697650535843 & 664,70855018579334 & 150104,02694537438 \\
S9 & 0,14824924 & 8,229083300000001 & 308,00322499999999 & 624000 & -15927018,300000001 & -9177,3234100000009 & 1 & 0 & 3568,7722176518409 & 49957,591485797835 & 2958,3249060838357 & 610,4473149407039 \\
S10 & 0,018557322099999999 & 0,42166340500000005 & 308,00322499999999 & 624000 & -8938691,6099999994 & -262,71059300000002 & 0 & 1 & 97789,376086757693 & 451490,35053931631 & 577,40901096931441 & 97211,967345223806 \\
S11 & 25,631088099999999 & 885,05505200000005 & 308,00322499999999 & 624000 & 8290,0620400000007 & -331,74889899999999 & 0 & 1 & 150768,73540924702 & 5249,6697945492215 & 664,70854977085423 & 150104,02660197637 \\
S12 & 25,631088099999999 & 885,05505200000005 & 100,879988 & 614000 & -217465,42600000001 & -1598,09349 & 0,050000000000000003 & 0,94999999999999996 & 289910,43926847569 & 5249,6697945492215 & 141137,64526446303 & 148772,79439416385 \\
S13 & 16,6404663 & 564,11606599999993 & 100,313575 & 613920 & -380571,96399999998 & -3174,27153 & 1 & 0 & 591473,30264416931 & 8916,6817999565319 & 445454,40532516805 & 146018,8973190012 \\
S14 & 16,6404663 & 564,11606599999993 & 91,795263000000006 & 313040 & -380571,96399999998 & -3162,4382700000001 & 0,90452481399999995 & 0,095475185800000001 & 588063,54843553877 & 8916,6817999565319 & 496574,01066940051 & 91489,537766138208 \\
S15 & 8,5410905400000008 & 304,89203700000002 & 96,712232700000001 & 610000 & -220829,09400000001 & -1778,6030800000001 & 0 & 1 & 291899,60091442842 & 25701,04262119202 & 138605,67739631992 & 153293,9235181085 \\
S16 & 8,5410905400000008 & 304,89203700000002 & 306,84578699999997 & 600000 & 7518,6372700000002 & -502,353161 & 0 & 1 & 152495,91886424375 & 25701,04262119202 & 611,19428316000494 & 151884,72407880597 \\
S17 & 0,44953124300000002 & 16,046949299999998 & 96,712232700000001 & 610000 & -390412,95899999997 & -3532,0898900000002 & 1 & 0 & 627583,94081187353 & 25700,386969543739 & 474290,07287041884 & 153293,86816390869 \\
S18 & 0,44953124300000002 & 16,046949299999998 & 88,413802399999994 & 310000 & -390412,95899999997 & -3518,7934399999999 & 0,90035230899999996 & 0,099647691100000005 & 623752,56529166317 & 25700,386969543739 & 528213,38604934292 & 95539,180354589946 \\
S19 & 4,8821120200000001 & 168,58191500000001 & 308,00322499999999 & 624000 & 8290,0620400000007 & -331,74889899999999 & 0 & 1 & 150768,73533106683 & 5249,6697730422002 & 664,70854964118575 & 150104,02669949387 \\
S20 & 4,8821120200000001 & 168,58191500000001 & 103,040773 & 614000 & -205871,71599999999 & -1483,5054299999999 & 0 & 1 & 268485,6030812664 & 5249,6697730422002 & 119712,80904775305 & 148772,79423834279 \\
S21 & 4,8821120200000001 & 168,58191500000001 & 92,539984200000006 & 311600 & -219799,87100000001 & -1456,9383 & 0,040048279899999997 & 0,95995171999999995 & 246902,12659233488 & 5249,6697730422002 & 154080,584369713 & 92821,543246768837 \\
S22 & 14,6033372 & 485,70152999999999 & 93,060852699999998 & 313840 & -394103,44799999997 & -3271,0736000000002 & 1 & 0 & 610469,47269011906 & 18675,879647564394 & 520517,73343972355 & 89951,738565620471 \\
S23 & 14,6033372 & 485,70152999999999 & 83,0483464 & 120000 & -394103,44799999997 & -3256,1666399999999 & 0,90381779799999995 & 0,096182201600000003 & 606174,03191922465 & 18675,879647564394 & 592706,81498746737 & 13467,217137189709 \\
S24 & 0,14603337199999999 & 4,8570152999999996 & 83,0483464 & 120000 & -394103,44799999997 & -3256,1666399999999 & 0,90381779799999995 & 0,096182201600000003 & 606174,03191922465 & 18675,879647564394 & 592706,81498746737 & 13467,217137189709 \\
S25 & 0,14603337199999999 & 4,8570152999999996 & 82,251933899999997 & 110000 & -394103,44799999997 & -3253,7647400000001 & 0,89680535400000005 & 0,103194646 & 605481,9271036213 & 18675,879647564394 & 598941,69738133554 & 6540,2296606559221 \\
S26 & 14,4573038 & 480,844514 & 83,0483464 & 120000 & -394103,44799999997 & -3256,1666399999999 & 0,90381779799999995 & 0,096182201600000003 & 606174,03294796916 & 18675,879661600524 & 592706,81577570504 & 13467,217172264167 \\
S27 & 14,4573038 & 480,844514 & 86,521558499999998 & 120000 & -211061,07699999999 & -1096,7272800000001 & 0,023204379399999999 & 0,97679562099999995 & 166973,95194808039 & 18675,879661600524 & 153506,73477581621 & 13467,217172264167 \\
S28 & 14,4573038 & 480,844514 & 306,84578699999997 & 110000 & 8112,3291200000003 & 218,712153 & 0 & 1 & 7103,4866127666219 & 18675,879661600524 & 563,25691239883884 & 6540,2296727001058 \\
S29 & 7,3687723900000002 & 263,04339999999996 & 88,413843499999999 & 310000 & -224177,81099999999 & -1638,5916199999999 & 0 & 1 & 248209,89483704217 & 25698,728224661583 & 152670,7883563764 & 95539,107159204854 \\
S30 & 7,3687723900000002 & 263,04339999999996 & 306,84578699999997 & 300000 & 8276,3640599999999 & -294,441732 & 0 & 1 & 93347,281147328264 & 25698,728224661583 & 607,90679544927571 & 92739,374326094505 \\
S31 & 7,3687723900000002 & 263,04339999999996 & 388,83178199999998 & 610000 & 93551,131500000003 & -258,75215500000002 & 0 & 1 & 168338,09681560812 & 25698,728224661583 & 15044,34662012949 & 153293,74965264736 \\
S32 & 15,9098629 & 567,93543700000009 & 344,80136499999998 & 600000 & 47365,232499999998 & -379,91224199999999 & 0 & 1 & 157062,74187944134 & 25699,927118793712 & 5178,0964435589203 & 151884,64571872583 \\
SZ33 & 16,632650000000002 & 593,73712499999999 & 96,712232700000001 & 610000 & -220829,09400000001 & -1778,6030800000001 & 0 & 1 & 291899,60048458906 & 25701,042587921947 & 138605,67678632087 & 153293,92369826816 \\
SZ34 & 16,632650000000002 & 593,73712499999999 & 96,711881300000002 & 610000 & -390413,83100000001 & -3532,1030999999998 & 1 & 0 & 627585,89881949045 & 25701,042587921947 & 474291,97391876811 & 153293,92369826816 \\
SZ35 & 15,319466800000001 & 509,51973299999997 & 93,060852699999998 & 313840 & -394102,86499999999 & -3271,0673299999999 & 0,99999674299999997 & 3,25693592e-06 & 610468,25076183456 & 18675,879633095323 & 520516,51171044674 & 89951,739051387864 \\
SZ36 & 15,319466800000001 & 509,51973299999997 & 96,398110299999999 & 313840 & -209981,342 & -1326,4756299999999 & 0,051547579199999999 & 0,94845242100000005 & 234255,67461656043 & 18675,879633095323 & 144303,93556517255 & 89951,739051387864 \\
SZ37 & 14,7374711 & 526,08416999999997 & 88,413848000000002 & 310000 & -224177,80600000001 & -1638,5915600000001 & 0 & 1 & 248209,88385178242 & 25698,728189533143 & 152670,77673862237 & 95539,107113160004 \\
SZ38 & 14,7374711 & 526,08416999999997 & 88,413825000000003 & 310000 & -403740,44300000003 & -3669,5257099999999 & 0,972537863 & 0,027462137399999999 & 653860,92122684466 & 25698,728189533143 & 558321,81411368458 & 95539,107113160004 \\
SZ39 & 14,4573038 & 480,844514 & 83,0483464 & 120000 & -394104,63500000001 & -3256,1812199999999 & 0,90382416099999996 & 0,096175839099999993 & 606177,04872467299 & 18675,879661600524 & 592709,83155240875 & 13467,217172264167 \\
SZ40 & 14,4573038 & 480,844514 & 86,521538899999996 & 120000 & -211062,264 & -1096,741 & 0,023209824099999998 & 0,97679017599999995 & 166976,7187156986 & 18675,879661600524 & 153509,50154343439 & 13467,217172264167 \\
\hline
\end{longtable}
\end{landscape}


\centering
\begin{longtable}{lrrrrr}
\caption{Stoffliche Zusammensetzung der Prozessströme des Doppelkolonnenmodells} \\
\hline
Stream & $x_{N_2}$ & $x_{O_2}$ & $x_{CO_2}$ & $x_{Ar}$ & $x_{H_2O}$ \\
 & (-) & (-) & (-) & (-) & (-) \\
\hline
S1 & 0,77289536299999995 & 0,20739875599999999 & 0,00039599762399999998 & 0,0092099447399999999 & 0,010099939400000001 \\
S2 & 0,77289536299999995 & 0,20739875599999999 & 0,00039599762399999998 & 0,0092099447399999999 & 0,010099939400000001 \\
S3 & 0,77289536299999995 & 0,20739875599999999 & 0,00039599762399999998 & 0,0092099447399999999 & 0,010099939400000001 \\
S4 & 0,77289536299999995 & 0,20739875599999999 & 0,00039599762399999998 & 0,0092099447399999999 & 0,010099939400000001 \\
S5 & 0,77289536299999995 & 0,20739875599999999 & 0,00039599762399999998 & 0,0092099447399999999 & 0,010099939400000001 \\
S6 & 0,77473283900000001 & 0,20789181400000001 & 0,00039693898400000003 & 0,0092318400500000005 & 0,0077465673399999997 \\
S7 & 6,9420724399999998e-07 & 4,2599835999999997e-06 & 3,5161362199999999e-08 & 1,5873828799999999e-07 & 0,99999485200000005 \\
S8 & 0,78109368099999998 & 0,20959868300000001 & 0 & 0,0093076368500000003 & 0 \\
S9 & 0 & 0 & 0 & 0 & 1 \\
S10 & 0 & 0 & 1 & 0 & 0 \\
S11 & 0,78109368099999998 & 0,20959868300000001 & 0 & 0,0093076368500000003 & 0 \\
S12 & 0,78109368099999998 & 0,20959868300000001 & 0 & 0,0093076368500000003 & 0 \\
S13 & 0,623186032 & 0,36089992100000001 & 0 & 0,015914047099999998 & 0 \\
S14 & 0,623186032 & 0,36089992100000001 & 0 & 0,015914047099999998 & 0 \\
S15 & 0,99985058199999999 & 1,3983849600000001e-07 & 0 & 0,00014927787700000001 & 0 \\
S16 & 0,99985058199999999 & 1,3983849600000001e-07 & 0 & 0,00014927787700000001 & 0 \\
S17 & 0,99972122900000004 & 3,2507700599999999e-07 & 0 & 0,00027844583199999998 & 0 \\
S18 & 0,99972122900000004 & 3,2507700599999999e-07 & 0 & 0,00027844583199999998 & 0 \\
S19 & 0,78109368099999998 & 0,20959868300000001 & 0 & 0,0093076368500000003 & 0 \\
S20 & 0,78109368099999998 & 0,20959868300000001 & 0 & 0,0093076368500000003 & 0 \\
S21 & 0,78109368099999998 & 0,20959868300000001 & 0 & 0,0093076368500000003 & 0 \\
S22 & 0,66762704500000003 & 0,31831353000000001 & 0 & 0,0140594247 & 0 \\
S23 & 0,66762704500000003 & 0,31831353000000001 & 0 & 0,0140594247 & 0 \\
S24 & 0,66762704500000003 & 0,31831353000000001 & 0 & 0,0140594247 & 0 \\
S25 & 0,66762704500000003 & 0,31831353000000001 & 0 & 0,0140594247 & 0 \\
S26 & 0,66762704500000003 & 0,31831353000000001 & 0 & 0,0140594247 & 0 \\
S27 & 0,66762704500000003 & 0,31831353000000001 & 0 & 0,0140594247 & 0 \\
S28 & 0,66762704500000003 & 0,31831353000000001 & 0 & 0,0140594247 & 0 \\
S29 & 0,99999382400000003 & 4,7026693399999998e-09 & 0 & 6,1709334200000002e-06 & 0 \\
S30 & 0,99999382400000003 & 4,7026693399999998e-09 & 0 & 6,1709334200000002e-06 & 0 \\
S31 & 0,99999382400000003 & 4,7026693399999998e-09 & 0 & 6,1709334200000002e-06 & 0 \\
S32 & 0,99985351200000006 & 1,37074423e-07 & 0 & 0,000146350761 & 0 \\
SZ33 & 0,99985058199999999 & 1,3983849600000001e-07 & 0 & 0,00014927787700000001 & 0 \\
SZ34 & 0,99985058199999999 & 1,3983849600000001e-07 & 0 & 0,00014927787700000001 & 0 \\
SZ35 & 0,66762704500000003 & 0,31831353000000001 & 0 & 0,0140594247 & 0 \\
SZ36 & 0,66762704500000003 & 0,31831353000000001 & 0 & 0,0140594247 & 0 \\
SZ37 & 0,99999382400000003 & 4,7026693399999998e-09 & 0 & 6,1709334200000002e-06 & 0 \\
SZ38 & 0,99999382400000003 & 4,7026693399999998e-09 & 0 & 6,1709334200000002e-06 & 0 \\
SZ39 & 0,66762704500000003 & 0,31831353000000001 & 0 & 0,0140594247 & 0 \\
SZ40 & 0,66762704500000003 & 0,31831353000000001 & 0 & 0,0140594247 & 0 \\
\hline
\end{longtable}


\chapter{Exergieberechnungen}
\label{ch:a_exergie}

\begin{landscape}
\begin{table}[p]
\centering
\caption{Definition von Brennstoff- und Produktexergie der betrachteten Komponenten für das Einzelkolonnenmodell}
\label{tab:exdef_components_single}
\begin{tabular}{lll}
\toprule
Komponente & $\dot{E}_{F,k}$ & $\dot{E}_{P,k}$ \\
\midrule
LK1 &
$\dot{W}_{LK1}$ &
$\dot{E}^{PH}_{2}-\dot{E}^{PH}_{1}$ \\

LK2 &
$\dot{W}_{LK2}$ &
$\dot{E}^{PH}_{4}-\dot{E}^{PH}_{3}$ \\

\addlinespace
ZK1 (dissipativ) &
$\dot{E}^{PH}_{2}-\dot{E}^{PH}_{3}$ &-- \\

ZK2 (dissipativ) &
$\dot{E}^{PH}_{4}-\dot{E}^{PH}_{5}$ &-- \\

\addlinespace
GW1 (Flash, dissipativ) &
$\dot{E}_{D,GW1}=\dot{E}_{5}-\dot{E}_{6}-\dot{E}_{7}$ &
-- \\

GW2 (Adsorber, dissipativ) &
$\dot{E}_{D,GW2}=\dot{E}_{6}-\dot{E}_{8}-\dot{E}_{9}-\dot{E}_{10}$ &
-- \\

\addlinespace
MW &
$\dot{E}_{F,MW}=(\dot{E}_{16}-\dot{E}_{15})+(\dot{E}_{28}-\dot{E}_{27})+(\dot{E}_{30}-\dot{E}_{29})$ &
$\dot{E}_{P,MW}=(\dot{E}_{11}-\dot{E}_{12})+(\dot{E}_{19}-\dot{E}_{20})$ \\
-- \\

\addlinespace
KOL &
$???$ & $???$\\

\addlinespace
RC (Reboiler/Kondensator) &
${E}^{PH}_{35}-\dot{E}^{PH}_{36} + {E}^{M}_{33}-\dot{E}^{M}_{34}$ &
${E}^{T}_{34}-\dot{E}^{T}_{33}$ \\

\addlinespace
D1 (dissipativ) &
$\dot{E}^{PH}_{12}-\dot{E}^{PH}_{13}$ &-- \\

D2 (dissipativ) &
$\dot{E}^{PH}_{18}-\dot{E}^{PH}_{19}$ &-- \\

\addlinespace
T &
$\dot{E}^{PH}_{16}-\dot{E}^{PH}_{17}$ &
$\dot{W}_{T}$ \\
\bottomrule
\end{tabular}
\end{table}
\end{landscape}


\begin{landscape}

\begin{table}[p]
\centering
\caption{Definition von Brennstoff- und Produktexergie der betrachteten Komponenten für das Doppelkolonnenmodell}
\label{tab:exdef_components_doppel}
\begin{tabular}{lll}
\toprule
Komponente & $\dot{E}_{F,k}$ & $\dot{E}_{P,k}$ \\
\midrule
LK1 &
$\dot{W}_{LK1}$ &
$\dot{E}^{PH}_{2}-\dot{E}^{PH}_{1}$ \\

LK2 &
$\dot{W}_{LK2}$ &
$\dot{E}^{PH}_{4}-\dot{E}^{PH}_{3}$ \\

\addlinespace
ZK1 (dissipativ) &
$\dot{E}^{PH}_{2}-\dot{E}^{PH}_{3}$ &
-- \\

ZK2 (dissipativ) &
$\dot{E}^{PH}_{4}-\dot{E}^{PH}_{5}$ &
-- \\

\addlinespace
GW1 (Flash, dissipativ) &
$\dot{E}_{D,GW1}=\dot{E}_{5}-\dot{E}_{6}-\dot{E}_{7}$ &
-- \\

GW2 (Adsorber, dissipativ) &
$\dot{E}_{D,GW2}=\dot{E}_{6}-\dot{E}_{8}-\dot{E}_{9}-\dot{E}_{10}$ &
-- \\

\addlinespace
MW &
$\dot{E}_{F,MW}=(\dot{E}_{16}-\dot{E}_{15})+(\dot{E}_{28}-\dot{E}_{27})+(\dot{E}_{30}-\dot{E}_{29})$ &
$\dot{E}_{P,MW}=(\dot{E}_{11}-\dot{E}_{12})+(\dot{E}_{19}-\dot{E}_{20})$ \\
-- \\

\addlinespace
KOLHP &
$\dot{E}_{F,KOLHP}=$ &
$\dot{E}_{P,KOLHP}=$ \\
-- \\

KOLLP &
$\dot{E}_{F,KOLLP}=$ &
$\dot{E}_{P,KOLLP}=$ \\
-- \\

\addlinespace
RC1 (Reboiler/Kondensator) &
$\dot{E}_{F,RC}=\dot{E}_{SZRC1}-\dot{E}_{SZRC2}$ &
$\dot{E}_{P,RC}=\dot{E}_{SZRC4}-\dot{E}_{SZRC3}$ \\
-- \\

RC2 (Reboiler/Kondensator) &
$\dot{E}_{F,RC}=$ &
$\dot{E}_{P,RC}=$ \\
-- \\

\addlinespace
D1 &
$\dot{E}_{F,D1}=$ &
$\dot{E}_{P,D1}=$ \\
-- \\

D2&
$\dot{E}_{F,D2}=$ &
$\dot{E}_{P,D2}=$ \\
-- \\

D3 &
$\dot{E}_{F,D3}=$ &
$\dot{E}_{P,D3}=$ \\
-- \\

D4&
$\dot{E}_{F,D4}=$ &
$\dot{E}_{P,D4}=$ \\
-- \\

\addlinespace
T &
$\dot{E}_{F,T}=$ &
$\dot{E}_{P,T}=$ \\
-- \\

\addlinespace
PK1 &
$\dot{W}_{PK1}$ &
$\dot{E}^{PH}_{31}-\dot{E}^{PH}_{30}$ \\
\bottomrule
\end{tabular}
\end{table}
\end{landscape}

\begin{table}[p]
\centering
\caption{Definition von Brennstoff- und Produktexergie ausgewählter Funktionsgruppen der Verdichtungsstrecke (Einzelkolonnenmodell)}
\label{tab:exdef_groups_single_total}
\begin{tabular}{lll}
\toprule
Funktionsgruppe & $\dot{E}_{F,g}$ & $\dot{E}_{P,g}$ \\
\midrule
ZK1 + LK2 &
$\dot{W}_{LK2}$ &
$\dot{E}_{4}-\dot{E}_{2}$ \\

LK1 + ZK1 + LK2 &
$\dot{W}_{LK1}+\dot{W}_{LK2}$ &
$\dot{E}_{4}-\dot{E}_{1}$ \\
\bottomrule
\end{tabular}
\end{table}

\begin{table}[p]
\centering
\caption{Berechnete exergetische Kenngrößen der betrachteten Anlagenkomponenten}
\label{tab:exergy-components}
\begin{tabular}{llrrrrrr}
\hline
Component & Type & $\dot{E}_F$ & $\dot{E}_P$ & $\dot{E}_D$ & $\dot{E}_L$ & $\varepsilon$ & $y_{D,k}$ \\
 &  & (W) & (W) & (W) & (W) & (-) & (-) \\
\hline
LK1 & Compressor & 3.33712e+06 & 2.93889e+06 & 398233 & 0 & 0.880666 & 0.0608209 \\\\
LK2 & Compressor & 3.5683e+06 & 3.16997e+06 & 398326 & 0 & 0.888371 & 0.0608351 \\\\
PK1 & Compressor & 607700 & 534420 & 73280.1 & 0 & 0.879414 & 0.0111918 \\\\
GW1 & Flash2 & 42460.2 & 4725.72 & 35285.5 & 2449.05 & 0.111297 & 0.00538904 \\\\
MW & MHeatX & 5.2411e+06 & 4.66321e+06 & 577891 & 0 & 0.889739 & 0.0882594 \\\\
RC & MHeatX & 5.1909e+06 & 5.02843e+06 & 162475 & 0 & 0.9687 & 0.0248143 \\\\
RC2 & MHeatX & 5.69621e+06 & 5.47171e+06 & 224507 & 0 & 0.960587 & 0.0342882 \\\\
MIX1 & Mixer & 82153.7 & 38354.4 & 43799.3 & 0 & 0.466862 & 0.00668932 \\\\
KOLHP & RadFrac & 5.2497e+06 & 5.15659e+06 & 93117.4 & 0 & 0.982262 & 0.0142215 \\\\
KOLLP & RadFrac & 2.2108e+06 & 1.77318e+06 & 437618 & 0 & 0.802055 & 0.0668359 \\\\
GW2 & Sep & 75398 & 38519.8 & 18653.9 & 18224.3 & 0.510886 & 0.00284894 \\\\
ZK1 & SimpleHeatExchanger & 579569 & 0 & 579569 & 0 & - & 0.0885158 \\\\
ZK2 & SimpleHeatExchanger & 784130 & 0 & 784130 & 0 & - & 0.119758 \\\\
T & Turbine & 274606 & 236350 & 38255.5 & 0 & 0.860689 & 0.00584264 \\\\
D1 & Valve & 894787 & 838903 & 55883.5 & 0 & 0.937545 & 0.0085349 \\\\
D2 & Valve & 26974.7 & 25185.3 & 1789.46 & 0 & 0.933662 & 0.000273298 \\\\
D3 & Valve & 952026 & 904480 & 47546.6 & 0 & 0.950057 & 0.00726165 \\\\
D4 & Valve & 3682.59 & 3323.86 & 358.726 & 0 & 0.902589 & 5.47871e-05 \\\\
System Discharge & Boundary & 0 & 0 & 0 & 669810 & - & - \\\\
\hline
\end{tabular}

\textbf{Exergetischer Global-Check}\\
\begin{tabular}{lll}
\hline
Größe & Formel & Wert \\
\hline
$\dot E_{tot,S7}$ & $\dot m_{S7}(e_{PH,S7}+e_{CH,S7})$ & 2449.05 \\
$\dot E_{tot,S9}$ & $\dot m_{S9}(e_{PH,S9}+e_{CH,S9})$ & 7977.23 \\
$\dot E_{tot,S10}$ & $\dot m_{S10}(e_{PH,S10}+e_{CH,S10})$ & 10247.1 \\
$\dot E_{tot,S25}$ & $\dot m_{S25}(e_{PH,S25}+e_{CH,S25})$ & 91491.7 \\
$\dot E_{tot,S28}$ & $\dot m_{S28}(e_{PH,S28}+e_{CH,S28})$ & 557645 \\
$\sum \dot E_{tot,i}$ (Discharge) & $\sum_{i \in \{S7,S9,S10,S25,S28\}} \dot E_{tot,i}$ & 669810 \\
\hline
Total Input ($E_{in}$) & $E_{F,\mathrm{LK1}} + E_{F,\mathrm{LK2}} + E_{F,\mathrm{PK1}} + \dot E_{S1}$ & 7.63175e+06 \\
Total Output ($E_{out}$) & $\dot E_{S32} + E_{P,\mathrm{T}}$ & 3.15653e+06 \\
$\sum E_D + \sum E_L$ & $\sum_k \dot E_{D,k} + \sum_k \dot E_{L,k}$ & 4.6612e+06 \\
Bilanz-Differenz [W] & $E_{in} - (E_{out} + \sum E_D + \sum E_L)$ & -185985 \\
Bilanz-Differenz [\%] & $100 \cdot \Delta E_{Bilanz}/E_{in}$ & -2.43699 \\
\hline
\end{tabular}
\end{table}


\chapter{Wirtschaftlichkeitsanalyse}
\label{ch:a_kosten}
